\FigureLayoutDeclare{img/2008/syllabus_paper_1_liberal/q18_fig1.png}{5.0cm}{2bp 53bp 42bp 34bp}

\chapter{2008年大纲I卷（文）}

\section{单选题}

\begin{problem}
函数 \( y = \sqrt{1 - x} + \sqrt{x} \) 的定义域为（\quad）

\choices
    {\( \{x \mid x \leqslant 1\} \)}
    {\( \{x \mid x \geqslant 0\} \)}
    {\(\{x \mid x \geqslant 1 \text{或} x \leqslant 0\}\)}
    {\(\{x\mid 0\leqslant x\leqslant 1\}\)}
\end{problem}

\begin{answer}
D
\end{answer}

\begin{problem}
汽车经过启动，加速行驶，匀速行驶，减速行驶之后停车，若把这一过程中汽车的行驶路程 \(s\) 看作时间 \(t\) 的函数，其图象可能是（\quad）

\choices
    {\choicebitmap[width=0.15\paperwidth]{img/2008/syllabus_paper_1_liberal/q2_optA.png}}
    {\choicebitmap[width=0.15\paperwidth]{img/2008/syllabus_paper_1_liberal/q2_optB.png}}
    {\choicebitmap[width=0.15\paperwidth]{img/2008/syllabus_paper_1_liberal/q2_optC.png}}
    {\choicebitmap[width=0.15\paperwidth]{img/2008/syllabus_paper_1_liberal/q2_optD.png}}
\end{problem}

\begin{answer}
A
\end{answer}

\begin{problem}
\(\left(1 + \frac{x}{2}\right)^5\) 的展开式中 \(x^2\) 的系数为（\quad）

\choices
    {\(10\)}
    {\(5\)}
    {\(\frac{5}{2}\)}
    {\(1\)}
\end{problem}

\begin{answer}
C
\end{answer}

\begin{problem}
曲线 \( y = x^{3} - 2x + 4 \) 在点 \((1, 3)\) 处的切线的倾斜角为（\quad）

\choices
    {\(30^{\circ}\)}
    {\(45^{\circ}\)}
    {\(60^{\circ}\)}
    {\(120^{\circ}\)}
\end{problem}

\begin{answer}
B
\end{answer}

\begin{problem}
在 \(\triangle ABC\) 中，\(\overrightarrow{AB} = \bs{c}\)，\(\overrightarrow{AC} = \bs{b}\). 若点 \(D\) 满足 \(\overrightarrow{BD} = 2\overrightarrow{DC}\)，则 \(\overrightarrow{AD} =\)（\quad）

\choices
    {\(\frac{2}{3} \bs{b} + \frac{1}{3} \bs{c}\)}
    {\(\frac{5}{3} \bs{c} - \frac{2}{3} \bs{b}\)}
    {\(\frac{2}{3} \bs{b} - \frac{1}{3} \bs{c}\)}
    {\(\frac{1}{3} \bs{b} + \frac{2}{3} \bs{c}\)}
\end{problem}

\begin{answer}
A
\end{answer}

\begin{problem}
\(y = (\sin x - \cos x)^{2} - 1\) 是（\quad）

\choices
    {最小正周期为 \(2\pi\) 的偶函数}
    {最小正周期为 \(2\pi\) 的奇函数}
    {最小正周期为 \(\pi\) 的偶函数}
    {最小正周期为 \(\pi\) 的奇函数}
\end{problem}

\begin{answer}
D
\end{answer}

\begin{problem}
已知等比数列 \(\{a_{n}\}\) 满足 \(a_{1} + a_{2} = 3\)，\(a_{2} + a_{3} = 6\)，则 \(a_{7} =\)（\quad）

\choices
    {\(64\)}
    {\(81\)}
    {\(128\)}
    {\(243\)}
\end{problem}

\begin{answer}
A
\end{answer}

\begin{problem}
若函数 \( y = f(x) \) 的图象与函数 \( y = \ln \sqrt{x} + 1 \) 的图象关于直线 \( y = x \) 对称，则 \( f(x) = \)（\quad）

\choices
    {\(\e^{2x - 2}\)}
    {\(\e^{2x}\)}
    {\(\e^{2x + 1}\)}
    {\(\e^{2x + 2}\)}
\end{problem}

\begin{answer}
A
\end{answer}

\begin{problem}
为得到函数 \( y = \cos \left( x + \frac{\pi}{3} \right) \) 的图象，只需将函数 \( y = \sin x \) 的图象（\quad）

\choices
    {向左平移 \(\frac{\pi}{6}\) 个长度单位}
    {向右平移 \(\frac{\pi}{6}\) 个长度单位}
    {向左平移 \(\frac{5\pi}{6}\) 个长度单位}
    {向右平移 \(\frac{5\pi}{6}\) 个长度单位}
\end{problem}

\begin{answer}
C
\end{answer}

\begin{problem}
若直线 \(\frac{x}{a} + \frac{y}{b} = 1\) 与圆 \(x^{2} + y^{2} = 1\) 有公共点，则（\quad）

\choices
    {\(a^2 + b^2 \leqslant 1\)}
    {\(a^2 + b^2 \geqslant 1\)}
    {\(\frac{1}{a^2} +\frac{1}{b^2}\leqslant 1\)}
    {\(\frac{1}{a^2} +\frac{1}{b^2}\geqslant 1\)}
\end{problem}

\begin{answer}
D
\end{answer}

\begin{problem}
已知三棱柱 \(ABC - A_{1}B_{1}C_{1}\) 的侧棱与底面边长都相等，\(A_{1}\) 在底面 \(ABC\) 内的射影为 \(\triangle ABC\) 的中心，则 \(AB_{1}\) 与底面 \(ABC\) 所成角的正弦值等于（\quad）

\choices
    {\(\frac{1}{3}\)}
    {\(\frac{\sqrt{2}}{3}\)}
    {\(\frac{\sqrt{3}}{3}\)}
    {\(\frac{2}{3}\)}
\end{problem}

\begin{answer}
B
\end{answer}

\begin{problem}
将 \(1\)，\(2\)，\(3\) 填入 \(3 \times 3\) 的方格中，要求每行，每列都没有重复数字，下面是一种填法，则不同的填写方法共有（\quad）

\begin{center}
\begin{tabular}{|c|c|c|}
\hline
\(1\) & \(2\) & \(3\) \\
\hline
\(3\) & \(1\) & \(2\) \\
\hline
\(2\) & \(3\) & \(1\) \\
\hline
\end{tabular}
\end{center}

\choices
    {\(6\text{ 种}\)}
    {\(12\text{ 种}\)}
    {\(24\text{ 种}\)}
    {\(48\text{ 种}\)}
\end{problem}

\begin{answer}
B
\end{answer}

\section{填空题}

\begin{problem}
若 \(x\)，\(y\) 满足约束条件 \(\begin{cases}
x + y \geqslant 0,\\
x - y + 3 \geqslant 0,\\
0 \leqslant x \leqslant 3,
\end{cases}\) 则 \(z = 2x - y\) 的最大值为\fillinblank{}.
\end{problem}

\begin{answer}
\(9\)
\end{answer}

\begin{problem}
已知抛物线 \( y = ax^2 - 1 \) 的焦点是坐标原点，则以抛物线与两坐标轴的三个交点为顶点的三角形面积为\fillinblank{}.
\end{problem}

\begin{answer}
\(2\)
\end{answer}

\begin{problem}
在 \(\triangle ABC\) 中，\(\angle A = 90^{\circ}\)，\(\tan B = \frac{3}{4}\)，若以 \(A\)，\(B\) 为焦点的椭圆经过点 \(C\)，则该椭圆的离心率 \(e = \)\fillinblank{}.
\end{problem}

\begin{answer}
\(\frac{1}{2}\)
\end{answer}

\begin{problem}
已知菱形 \(ABCD\) 中，\(AB = 2\)，\(\angle A = 120^{\circ}\)，沿对角线 \(BD\) 将 \(\triangle ABD\) 折起，使二面角 \(A - BD - C\) 为 \(120^{\circ}\)，则点 \(A\) 到 \(\triangle BCD\) 所在平面的距离等于\fillinblank{}.
\end{problem}

\begin{answer}
\(\frac{\sqrt{3}}{2}\)
\end{answer}

\section{解答题}

\begin{problem}
\begin{enumerate}
\item 设 \(\triangle ABC\) 的内角 \(A\)，\(B\)，\(C\) 所对的边长分别为 \(a\)，\(b\)，\(c\)，且 \(a \cos B = 3\)，\(b \sin A = 4\). 求边长 \(a\)；
\item 若 \(\triangle ABC\) 的面积 \(S = 10\)，求 \(\triangle ABC\) 的周长 \(l\).
\end{enumerate}
\end{problem}

\begin{solution}
由正弦定理，\(\dfrac{a}{\sin A}=\dfrac{b}{\sin B}\)，所以
\(
a\sin B=b\sin A=4
\).
于是
\(
\frac{\cos B}{\sin B}=\frac{a\cos B}{a\sin B}=\frac{3}{4}
\).
由于 \(a\cos B=3>0\)，故 \(B\) 为锐角，从而
\(
\sin B=\frac{4}{5}
\)，\(\cos B=\frac{3}{5}\).
因此
\(
a=\frac{4}{\sin B}=5
\).

当 \(S=10\) 时，利用三角形面积公式，得
\(
10=\frac{1}{2}bc\sin A=\frac{1}{2}c\left(b\sin A\right)=2c
\).
所以 \(c=5\). 再由余弦定理，得
\(
b^2=a^2+c^2-2ac\cos B=25+25-2\times5\times5\times\frac{3}{5}=20
\).
因为 \(b>0\)，所以 \(b=2\sqrt{5}\)，三角形的周长为
\(
l=a+b+c=5+2\sqrt{5}+5=10+2\sqrt{5}
\).
\end{solution}

\begin{problem}
四棱锥 \(A - BCDE\) 中，底面 \(BCDE\) 为矩形，侧面 \(ABC \perp\) 底面 \(BCDE\)，\(BC = 2\)，\(CD = \sqrt{2}\)，\(AB = AC\).

\begin{enumerate}
\item 证明： \(AD \perp CE\)；
\item 设侧面 \(ABC\) 为等边三角形，求二面角 \(C - AD - E\) 的大小.
\end{enumerate}

\bitmapfigure[max width=0.42\linewidth]{img/2008/syllabus_paper_1_liberal/q18_fig1.png}
\end{problem}

\begin{solution}
建立空间直角坐标系，令 \(B(0,0,0)\)，\(C(2,0,0)\)，\(D(2,\sqrt{2},0)\)，\(E(0,\sqrt{2},0)\). 底面在平面 \(z=0\) 内，且平面 \(ABC\) 与底面垂直并包含 \(BC\)，所以可令平面 \(ABC\) 为 \(y=0\). 由 \(AB=AC\)，点 \(A\) 在 \(BC\) 的垂直平分面内，故设 \(A(1,0,h)\)，其中 \(h>0\).

于是 \(\overrightarrow{AD}=(1,\sqrt{2},-h)\)，\(\overrightarrow{CE}=(-2,\sqrt{2},0)\)，从而
\(
\overrightarrow{AD}\cdot\overrightarrow{CE}=1\times(-2)+\sqrt{2}\times\sqrt{2}+(-h)\times0=0
\).
所以 \(AD\perp CE\).

当侧面 \(ABC\) 为等边三角形时，\(BC=2\)，故等边三角形的高为 \(\sqrt{3}\)，即 \(h=\sqrt{3}\)，从而 \(A(1,0,\sqrt{3})\).

设 \(\theta\) 为二面角 \(C-AD-E\) 的大小，并取 \(\bs{d}=\overrightarrow{DA}=(-1,-\sqrt{2},\sqrt{3})\)，\(\bs{u}=\overrightarrow{DC}=(0,-\sqrt{2},0)\)，\(\bs{v}=\overrightarrow{DE}=(-2,0,0)\). 将 \(\bs{u}\) 和 \(\bs{v}\) 分别投影到垂直于 \(AD\) 的平面内，所得向量就是两个面的平面角的两条边.

直接计算得 \(\lvert\bs{d}\rvert^2=6\)，\(\bs{u}\cdot\bs{d}=2\)，\(\bs{v}\cdot\bs{d}=2\)，\(\bs{u}\cdot\bs{v}=0\). 因此
\(
\bs{u}_{\perp}=\bs{u}-\frac{\bs{u}\cdot\bs{d}}{\lvert\bs{d}\rvert^2}\bs{d}=\bs{u}-\frac{1}{3}\bs{d}=\left(\frac{1}{3},-\frac{2\sqrt{2}}{3},-\frac{\sqrt{3}}{3}\right)
\)，
\(
\bs{v}_{\perp}=\bs{v}-\frac{\bs{v}\cdot\bs{d}}{\lvert\bs{d}\rvert^2}\bs{d}=\bs{v}-\frac{1}{3}\bs{d}=\left(-\frac{5}{3},\frac{\sqrt{2}}{3},-\frac{\sqrt{3}}{3}\right)
\).
从而
\(
\bs{u}_{\perp}\cdot\bs{v}_{\perp}=-\frac{2}{3}
\)，\(\lvert\bs{u}_{\perp}\rvert^2=\frac{4}{3}\)，\(\lvert\bs{v}_{\perp}\rvert^2=\frac{10}{3}\).
于是
\(
\cos\theta=\frac{\bs{u}_{\perp}\cdot\bs{v}_{\perp}}{\lvert\bs{u}_{\perp}\rvert\lvert\bs{v}_{\perp}\rvert}=\frac{-\frac{2}{3}}{\sqrt{\frac{4}{3}}\sqrt{\frac{10}{3}}}=-\frac{1}{\sqrt{10}}
\).
故 \(\theta=\arccos\left(-\frac{1}{\sqrt{10}}\right)\).
\end{solution}

\begin{problem}
\begin{enumerate}
\item 在数列 \(\{a_{n}\}\) 中，\(a_{1} = 1\)，\(a_{n+1} = 2a_{n} + 2^{n}\). 设 \( b_{n}=\frac{a_{n}}{2^{n-1}} \). 证明： 数列 \( \{b_{n}\} \) 是等差数列；
\item 求数列 \(\{a_{n}\}\) 的前 \(n\) 项和 \( S_{n} \).
\end{enumerate}
\end{problem}

\begin{solution}
由递推式，
\(
b_{n+1}=\frac{a_{n+1}}{2^n}=\frac{2a_n+2^n}{2^n}=\frac{a_n}{2^{n-1}}+1=b_n+1
\).
所以 \(b_{n+1}-b_n=1\)，数列 \(\{b_n\}\) 是公差为 \(1\) 的等差数列. 又 \(b_1=a_1=1\)，故
\(
b_n=b_1+(n-1)\times1=n
\).
从而
\(
a_n=2^{n-1}b_n=n2^{n-1}
\).

因此
\(
S_n=\sum_{k=1}^{n}k2^{k-1}
\).
将上式乘以 \(2\)，并作换指标变形，得
\(
2S_n=\sum_{k=1}^{n}k2^k=\sum_{k=2}^{n+1}(k-1)2^{k-1}=\sum_{k=2}^{n}(k-1)2^{k-1}+n2^n
\).
于是
\(
S_n=2S_n-S_n=n2^n-1-\sum_{k=2}^{n}2^{k-1}=n2^n-1-(2^n-2)=(n-1)2^n+1
\).
故数列 \(\{a_n\}\) 的前 \(n\) 项和为 \(S_n=(n-1)2^n+1\).
\end{solution}

\begin{problem}
已知 \(5\text{ 只}\)动物中有 \(1\text{ 只}\)患有某种疾病，需要通过化验血液来确定患病的动物. 血液化验结果呈阳性的即为患病动物，呈阴性即没患病，下面是两种化验方案：

\begin{center}
\begin{tabular}{|c|p{0.68\linewidth}|}
\hline
方案甲： & 逐个化验，直到能确定患病动物为止. \\
\hline
方案乙： & 先任取 \(3\text{ 只}\)，将它们的血液混在一起化验. 若结果呈阳性则表明患病动物为这 \(3\text{ 只}\)中的 \(1\text{ 只}\)，然后再逐个化验，直到能确定患病动物为止；若结果呈阴性则在另外 \(2\text{ 只}\)中任取 \(1\text{ 只}\)化验. \\
\hline
\end{tabular}
\end{center}

求依方案甲所需化验次数不少于依方案乙所需化验次数的概率.
\end{problem}

\begin{solution}
将动物按方案甲的化验顺序编号为1，2，3，4，5，并设患病动物等可能地是这五只中的任意一只. 方案甲中，若第1，2，3，4只患病，分别需要化验1，2，3，4次；若前四只均未患病，则第5只必患病，不必再化验第5只. 因而方案甲所需次数依次为
\(
N_{\text{甲}}=(1,2,3,4,4)
\).

方案乙先把第1，2，3只混合化验. 若患病动物是第1只，混合化验呈阳性，再化验第1只即可确定，所需总次数为2次；若患病动物是第2或第3只，则混合化验后需要先后化验前两只，所需总次数均为3次. 若患病动物是第4或第5只，混合化验呈阴性，再任取另外两只中的一只化验：阳性时直接确定，阴性时另一只必患病，所需总次数均为2次. 所以
\(
N_{\text{乙}}=(2,3,3,2,2)
\).

比较两种方案可知，\(N_{\text{甲}}\geqslant N_{\text{乙}}\) 的情形是患病动物为第3，4，5只，共3种. 因此所求概率为
\(
P\left(N_{\text{甲}}\geqslant N_{\text{乙}}\right)=\frac{3}{5}
\).
\end{solution}

\begin{problem}
\begin{enumerate}
\item 已知函数 \(f(x) = x^{3} + ax^{2} + x + 1\)，\(a \in \R\). 讨论函数 \( f(x) \) 的单调区间；
\item 设函数 \( f(x) \) 在区间 \( \left(-\frac{2}{3}, -\frac{1}{3}\right) \) 内是减函数，求 \( a \) 的取值范围.
\end{enumerate}
\end{problem}

\begin{solution}
\begin{enumerate}
\item 对函数求导，得
\(
f'(x)=3x^2+2ax+1
\).
其判别式为
\(
\Delta=(2a)^2-4\times3\times1=4(a^2-3)
\).

当 \(a^2<3\) 时，\(\Delta<0\)，且二次项系数为正，所以 \(f'(x)>0\) 对任意 \(x\in\R\) 成立，函数在 \(\R\) 上单调递增.

当 \(a^2=3\) 时，\(f'(x)\geqslant0\)，且只有一个零点，函数仍在 \(\R\) 上单调递增.

当 \(a^2>3\) 时，令 \(x_1=\frac{-a-\sqrt{a^2-3}}{3}\)，\(x_2=\frac{-a+\sqrt{a^2-3}}{3}\)，则 \(x_1<x_2\)，并且
\(
f'(x)=3(x-x_1)(x-x_2)
\).
因此函数在 \(\left(-\infty,x_1\right)\) 和 \(\left(x_2,+\infty\right)\) 上单调递增，在 \(\left(x_1,x_2\right)\) 上单调递减.

\item 因为函数在开区间 \(\left(-\frac{2}{3},-\frac{1}{3}\right)\) 内为减函数，且 \(f\) 可导，所以该区间内有 \(f'(x)\leqslant0\). 由连续性，令 \(x\) 趋近于 \(-\frac{1}{3}\)，得到
\(
f'\left(-\frac{1}{3}\right)=3\times\frac{1}{9}-\frac{2a}{3}+1=\frac{4-2a}{3}\leqslant0
\).
所以 \(a\geqslant2\).

反之，若 \(a\geqslant2\)，由于区间内 \(x<0\)，有 \(2ax\leqslant4x\)，从而
\(
f'(x)=3x^2+2ax+1\leqslant3x^2+4x+1=(3x+1)(x+1)<0
\).
其中 \(-\frac{2}{3}<x< -\frac{1}{3}\). 因此 \(f(x)\) 在该区间内为减函数. 故 \(a\) 的取值范围是 \([2,+\infty)\).
\end{enumerate}
\end{solution}

\begin{problem}
\begin{enumerate}
\item 双曲线的中心为原点 \(O\)，焦点在 \(x\) 轴上，两条渐近线分别为 \(l_{1}\)，\(l_{2}\)，经过右焦点 \(F\) 垂直于 \( l_{1} \) 的直线分别交 \(l_{1}\)，\(l_{2}\) 于 \(A\)，\(B\) 两点. 已知 \( \left|\overrightarrow{OA}\right| \)， \( \left|\overrightarrow{AB}\right| \)， \( \left|\overrightarrow{OB}\right| \) 成等差数列，且 \( \overrightarrow{BF} \) 与 \( \overrightarrow{FA} \) 同向. 求双曲线的离心率；
\item 设 \(AB\) 被双曲线所截得的线段的长为 \(4\)，求双曲线的方程.
\end{enumerate}
\end{problem}

\begin{solution}
设双曲线方程为
\(
\frac{x^2}{a^2}-\frac{y^2}{b^2}=1
\).
其中 \(a>0\)，\(b>0\). 令 \(c=\sqrt{a^2+b^2}\)，则右焦点为 \(F(c,0)\)，离心率为 \(e=\frac{c}{a}\). 两条渐近线可取为
\(
l_1:y=\frac{b}{a}x
\)
，\(l_2:y=-\frac{b}{a}x.
\)
过 \(F\) 且垂直于 \(l_1\) 的直线方程为
\(
y=-\frac{a}{b}(x-c)
\).
分别与两条渐近线联立，得到 \(A\left(\frac{a^2}{c},\frac{ab}{c}\right)\) 和 \(B\left(\frac{a^2c}{a^2-b^2},-\frac{abc}{a^2-b^2}\right)\).

向量 \(\overrightarrow{FA}\) 的横坐标为负. 由 \(\overrightarrow{BF}\) 与 \(\overrightarrow{FA}\) 同向，点 \(B\) 必在 \(F\) 的右侧，从而 \(a>b\)，以下分母 \(a^2-b^2\) 为正. 由坐标计算
\(
\left|\overrightarrow{OA}\right|=\sqrt{\left(\frac{a^2}{c}\right)^2+\left(\frac{ab}{c}\right)^2}=a
\).
又有
\(
\left|\overrightarrow{FA}\right|=\sqrt{\left(-\frac{b^2}{c}\right)^2+\left(\frac{ab}{c}\right)^2}=b
\).
以及
\(
\left|\overrightarrow{FB}\right|=\frac{bc^2}{a^2-b^2}
\).
由于 \(F\) 位于 \(A\)，\(B\) 之间，故
\(
\left|\overrightarrow{AB}\right|=\left|\overrightarrow{AF}\right|+\left|\overrightarrow{FB}\right|=b+\frac{bc^2}{a^2-b^2}=\frac{2a^2b}{a^2-b^2}
\).
另外
\(
\left|\overrightarrow{OB}\right|=\sqrt{\left(\frac{a^2c}{a^2-b^2}\right)^2+\left(\frac{abc}{a^2-b^2}\right)^2}=\frac{ac^2}{a^2-b^2}
\).

由题意，\(\left|\overrightarrow{OA}\right|\)，\(\left|\overrightarrow{AB}\right|\)，\(\left|\overrightarrow{OB}\right|\) 依次成等差数列，所以
\(
2\left|\overrightarrow{AB}\right|=\left|\overrightarrow{OA}\right|+\left|\overrightarrow{OB}\right|
\).
代入上面的长度，得
\(
\frac{4a^2b}{a^2-b^2}=a+\frac{ac^2}{a^2-b^2}
\).
两边乘以 \(a^2-b^2\)，并利用 \(c^2=a^2+b^2\)，得
\(
4a^2b=a(a^2-b^2)+ac^2=2a^3
\).
所以 \(b=\frac{a}{2}\)，从而
\(
e=\frac{c}{a}=\sqrt{1+\frac{b^2}{a^2}}=\sqrt{1+\frac{1}{4}}=\frac{\sqrt{5}}{2}
\).

此时 \(c=\frac{\sqrt{5}}{2}a\)，且直线 \(AB\) 的斜率为 \(-2\)，故其方程为
\(
y=-2x+2c=-2x+\sqrt{5}a
\).
将其代入双曲线方程 \(x^2-4y^2=a^2\)，得到
\(
x^2-4\left(-2x+\sqrt{5}a\right)^2=a^2
\).
整理得
\(
15x^2-16\sqrt{5}ax+21a^2=0
\).
其两个根为
\(
x_1=\frac{7\sqrt{5}}{15}a
\)，\(x_2=\frac{3\sqrt{5}}{5}a\).
所以两交点的横坐标之差为
\(
x_2-x_1=\frac{2\sqrt{5}}{15}a
\).
直线 \(AB\) 的斜率为 \(-2\)，故双曲线截得的线段长为
\(
\sqrt{1+(-2)^2}\left(x_2-x_1\right)=\sqrt{5}\cdot\frac{2\sqrt{5}}{15}a=\frac{2a}{3}
\).
由题设该长度为 \(4\)，得 \(a=6\)，从而 \(b=3\). 所以双曲线方程为
\(
\frac{x^2}{36}-\frac{y^2}{9}=1
\).
\end{solution}
